% Generated from sacct COMPLETED jobs, 2026-08-26. Per object, single A40.
% Fit times: ours from per-epoch checkpoint timestamps over 53 rung27+MCFM(v2_D) runs,
% which excludes queue and requeue time. Baselines from job Elapsed.
\begin{table}[t]
\centering
\small
\setlength{\tabcolsep}{6pt}
\begin{tabular}{llcc}
\toprule
Method & Adaptation & Per-shape fit & Inference \\
\midrule
Frozen TRELLIS.2 \cite{trellis2} & none            & ---              & 15\,min \\
MeshNCA \cite{meshnca}           & per-shape       & 16\,min          & 2\,min  \\
L4GM \cite{l4gm}                 & feed-forward    & ---              & 3\,min  \\
DynaMesh (ours)                  & per-shape LoRA  & \textbf{2.8\,h}  & 15\,min \\
\midrule
SV4D 2.0 \cite{sv4d2}            & feed-forward    & ---              & 11\,min \\
DG4D \cite{dg4d}                 & per-shape       & 8\,min           & $<$1\,min \\
\bottomrule
\end{tabular}
\caption{\textbf{Cost per object.} Wall-clock on a single NVIDIA A40, for one object and a
150-frame sequence. \emph{Per-shape fit} is optimisation performed for that object before
any output can be produced; a dash means the method is feed-forward and fits nothing.
\emph{Inference} is generating the full sequence once fitting is complete. Ours fits a
rank-4 LoRA for 30 epochs --- 2.33M trainable parameters, 0.18\% of the 1.3B texture flow
--- median 2.8\,h (IQR 2.0--3.4) over 53 runs, measured from per-epoch checkpoint
timestamps so queue and requeue time are excluded. Baseline figures are job wall-clock
over 21--27 runs each. Our inference is comparable to the feed-forward baselines; the
additional cost is the one-time adaptation, and it is what buys a mesh whose geometry is
provably unchanged across the sequence. The temporal component is parameter-free and adds
no measurable cost to either column, since the blend is applied to cached conditioning
tokens before the flow runs. Timings are under normal cluster contention on mixed
a40/L40S nodes, not isolated benchmarks.}
\label{tab:timing}
\end{table}
